% !PRINTED 521
\weberpage{521}

Aus diesem Satze ziehen wir hier eine wichtige Folgerung: Die Functionen

\weq{1}{\varphi = \varphi_0 t^h + \varphi_1 t^{h-1} + \cdots + \varphi_h,}

in denen die Co\"efficienten $\varphi_0, \varphi_1, \ldots, \varphi_h$ ganze oder gebrochene Functionale in $\Omega$ sind, aber von den Variablen $t$ frei angenommen werden, geh\"oren selbst zu den Functionalen im K\"orper $\Omega$. Von ihnen gilt der Satz:

2. \gesperrt{Ein Functional $\varphi$ ist nur dann ganz, wenn die Co\"efficienten $\varphi_0, \varphi_1, \ldots, \varphi_h$ ganze Functionale sind.}

Um ihn zu beweisen, nehmen wir an, es seien $\varphi_0, \varphi_1, \ldots, \varphi_h$ nicht alle zugleich ganz. Bestimmen wir ihren Hauptnenner $\mu$ und setzen

\[
\mu\varphi_0 = \alpha_0, \quad \mu\varphi_1 = \alpha_1, \quad \ldots, \quad \mu\varphi_h = \alpha_h,
\]

so sind $\alpha_0, \alpha_1, \ldots, \alpha_h$ ganze Functionale, und $\mu$ enth\"alt wenigstens einen Primfactor $\pi$, der nicht zugleich in allen Z\"ahlern $\alpha_0,, \alpha_1, \ldots, \alpha_h$ aufgeht % [sic] doubled comma in print
(vgl. den Schluss des vor. Paragraphen).

Dann ist die Function

\weq{2}{\chi = \mu\varphi = \alpha_0 t^h + \alpha_1 t^{h-1} + \cdots + \alpha_h,}

deren Co\"efficienten nun ganz sind, gewiss ein ganzes Functional.

Nehmen wir nun an, es sei $\varphi$ selbst ganz, so ist $\mu\varphi$ durch $\pi$ theilbar, w\"ahrend doch nicht s\"ammtliche Co\"efficienten $\alpha_0, \alpha_1, \ldots, \alpha_h$ durch $\pi$ theilbar sind. Wenn nun $\varphi$ ein ganzes Functional ist, so gen\"ugt es einer Gleichung von der Form

\weq{3}{\varphi^m = C_1 \varphi^{m-1} + C_2 \varphi^{m-2} + \cdots + C_m,}

in der die Co\"efficienten $C_1, C_2, \ldots, C_m$ ganze rationale Functionale sind.

Diese Functionale setzen wir nach \S.\,136, (2) in die Form

\[
C_1 = \frac{a_1 E_1}{E}, \quad C_2 = \frac{a_2 E_2}{E}, \quad \ldots, \quad C_m = \frac{a_m E_m}{E},
\]

worin $a_1, a_2, \ldots, a_m$ die absoluten Werthe von $C_1, C_2, \ldots, C_m$, also nat\"urliche ganze Zahlen (oder Null), und $E, E_1, \ldots, E_m$ primitive ganze Functionen, also Einheiten, sind.

Dann ergiebt die Gleichung (3):

\[
E\varphi^m = a_1 E_1 \varphi^{m-1} + a_2 E_2 \varphi^{m-2} + \cdots + a_m E_m,
\]

und durch Multiplication mit $\mu^m$:

\weq{4}{E\chi^m = \mu\left(a_1 E_1 \chi^{m-1} + a_2 E_2 \mu \chi^{m-2} + \cdots + a_m E_m \mu^{m-1}\right).}
